\section{Simplicial and Singular Homology}

For a CW complex, the group $\pi_n(X)$ depends only on the $(n+1)$-skeleton.  Homology groups provide a more computable invariant with the same skeletal dependence; for spheres, $H_i(S^n)$ is zero for $i>n$ and agrees with $\pi_i(S^n)$ for $1\leq i\leq n$.

The graph $X_1$ with vertices $x,y$ and four oriented edges $a,b,c,d$ has $C_1=\mathbb{Z}\langle a,b,c,d\rangle$, $C_0=\mathbb{Z}\langle x,y\rangle$, and $\partial_1(a)=\partial_1(b)=\partial_1(c)=\partial_1(d)=y-x$.  Thus
\[
\partial_1(ka+\ell b+mc+nd)=(k+\ell+m+n)(y-x),
\]
and $\ker\partial_1$ has basis $a-b,b-c,c-d$.  Attaching a 2-cell $A$ along $a-b$ gives $\partial_2(A)=a-b$, $H_1(X_2)\cong\mathbb{Z}^2$ generated by $b-c,c-d$.  Attaching a second 2-cell $B$ along the same cycle gives $H_1(X_3)\cong\mathbb{Z}^2$ and $H_2(X_3)=\mathbb{Z}\langle A-B\rangle$.  Attaching a 3-cell $C$ along that sphere gives $\partial_3(C)=A-B$, killing this $H_2$ class.  This is the chain-complex pattern $H_n=\ker\partial_n/\operatorname{Im}\partial_{n+1}$.

\textbf{$\Delta$-Complexes}

An ordered $n$-simplex is the convex hull $[v_0,\ldots,v_n]$ of $n+1$ affinely independent points.  The standard simplex is
\[
\Delta^n=\{(t_0,\ldots,t_n)\in\mathbb{R}^{n+1}:\sum_i t_i=1,\ t_i\geq0\}.
\]
The ordering determines orientations and the canonical affine map $(t_i)\mapsto\sum_i t_i v_i$; the $t_i$ are barycentric coordinates.  A face is obtained by deleting one vertex, with the induced vertex ordering.  Write $\partial\Delta^n$ for the union of faces and $\dot\Delta^n=\Delta^n-\partial\Delta^n$.

A $\Delta$-complex structure on $X$ is a collection $\sigma_\alpha:\Delta^n\to X$ such that:
\begin{enumerate}
\item $\sigma_\alpha|_{\dot\Delta^n}$ is injective and these open images partition $X$;
\item the restriction to every face is one of the structure maps, using the induced ordering;
\item $A\subset X$ is open exactly when $\sigma_\alpha^{-1}(A)$ is open in every $\Delta^n$.
\end{enumerate}
Equivalently, the complex is obtained by gluing ordered simplices along ordered faces.  This gives a CW structure whose cells are the open simplices.  A simplicial complex is the special case in which simplices are determined by their vertex sets.

\section*{Simplicial Homology}

Let $\Delta_n(X)$ be the free abelian group on the open $n$-simplices, written as finite sums of characteristic maps.  Define
\[
\partial_n(\sigma_\alpha)=\sum_{i=0}^n(-1)^i\sigma_\alpha|[v_0,\ldots,\widehat v_i,\ldots,v_n].
\]

\begin{lemma}[Composition of boundary maps is zero]
\label{lem:boundary-composition-zero}
\dcref{ch2:2.1}
\group{Singular Homology}

$\partial_{n-1}\partial_n=0$ for every $n$.
\end{lemma}
\begin{proof}
\sketch In the double boundary sum, each codimension-two face occurs twice, according to the order in which its two deleted vertices are removed.  The two signs differ by a factor of $-1$.
\end{proof}

A chain complex is a sequence of abelian groups and homomorphisms with consecutive composite zero.  Its $n$-th homology group is
\[
H_n=\ker\partial_n/\operatorname{Im}\partial_{n+1}.
\]
Elements of the kernel are cycles, elements of the image are boundaries, and equal cosets are homologous.  For a $\Delta$-complex, write this group as $H_n^\Delta(X)$.

\begin{example}[Simplicial homology of $S^1$]
\label{ex:simplicial-homology-circle}
\uses{lem:boundary-composition-zero}
\dcref{ch2:2.2}
\group{Singular Homology}

For the one-vertex, one-edge structure on $S^1$, $\partial_1=0$ and $\Delta_n(S^1)=0$ for $n\geq2$.  Hence
\[
H_n^\Delta(S^1)\cong\begin{cases}\mathbb{Z}&n=0,1,\\0&n\geq2.\end{cases}
\]
\end{example}

\begin{example}[Simplicial homology of the torus]
\label{ex:simplicial-homology-torus}
\uses{lem:boundary-composition-zero}
\dcref{ch2:2.3}
\group{Singular Homology}

For the one-vertex, three-edge, two-triangle structure on $T$, $\partial_1=0$ and $\partial_2U=\partial_2L=a+b-c$.  Thus
\[
H_n^\Delta(T)\cong\begin{cases}\mathbb{Z}&n=0,2,\\\mathbb{Z}\oplus\mathbb{Z}&n=1,\\0&n\geq3.\end{cases}
\]
The classes of $a,b$ generate $H_1^\Delta(T)$ and $U-L$ generates $H_2^\Delta(T)$.
\end{example}

\begin{example}[Simplicial homology of $\mathbb{RP}^2$]
\label{ex:simplicial-homology-rp2}
\uses{lem:boundary-composition-zero}
\dcref{ch2:2.4}
\group{Singular Homology}

For the two-vertex, three-edge, two-triangle structure on $\mathbb{RP}^2$, $\operatorname{Im}\partial_1=\langle w-v\rangle$, $\partial_2U=-a+b+c$, and $\partial_2L=a-b+c$.  Therefore
\[
H_0^\Delta(\mathbb{RP}^2)\cong\mathbb{Z},\qquad H_1^\Delta(\mathbb{RP}^2)\cong\mathbb{Z}_2,\qquad H_2^\Delta(\mathbb{RP}^2)=0.
\]
\end{example}

\begin{example}[$\Delta$-complex structure on $S^n$]
\label{ex:delta-complex-sn}
\uses{lem:boundary-composition-zero}
\dcref{ch2:2.5}
\group{Singular Homology}

Glue two copies $U,L$ of $\Delta^n$ along their boundaries by the identity.  Then $\ker\partial_n=\mathbb{Z}\langle U-L\rangle$, so $H_n^\Delta(S^n)\cong\mathbb{Z}$ for this structure.
\end{example}

\section*{Singular Homology}

A singular $n$-simplex in $X$ is a continuous map $\sigma:\Delta^n\to X$.  Let $C_n(X)$ be the free abelian group on singular $n$-simplices and use the same alternating face formula for $\partial_n$.  The preceding lemma gives
\[
H_n(X)=\ker\partial_n/\operatorname{Im}\partial_{n+1}.
\]
The singular complex $S(X)$ has one ordered $n$-simplex for each singular simplex of $X$, with faces given by restrictions; consequently $H_n^\Delta(S(X))=H_n(X)$.

\begin{proposition}[Homology of disjoint union]
\label{prop:homology-disjoint-union}
\uses{lem:boundary-composition-zero}
\dcref{ch2:2.6}
\group{Singular Homology}

If $X=\coprod_\alpha X_\alpha$ is the decomposition into path-components, then
\[
H_n(X)\cong\bigoplus_\alpha H_n(X_\alpha).
\]
\end{proposition}
\begin{proof}
\sketch Every singular simplex has path-connected image, so the chain complex splits as the direct sum of the chain complexes of the components.
\end{proof}

\begin{proposition}[Zero-dimensional homology of path-connected space]
\label{prop:h0-path-connected}
\uses{prop:homology-disjoint-union}
\dcref{ch2:2.7}
\group{Singular Homology}

If $X$ is nonempty and path-connected, $H_0(X)\cong\mathbb{Z}$.  In general, $H_0(X)$ is a direct sum of one copy of $\mathbb{Z}$ for each path-component.
\end{proposition}
\begin{proof}
\sketch The augmentation $\varepsilon:C_0(X)\to\mathbb{Z}$ sums coefficients.  Path-connectedness supplies a path from a chosen basepoint to every point, showing $\ker\varepsilon=\operatorname{Im}\partial_1$.
\end{proof}

\begin{proposition}[Homology of a point]
\label{prop:homology-point}
\dcref{ch2:2.8}
\group{Singular Homology}

For a point, $H_0\cong\mathbb{Z}$ and $H_n=0$ for $n>0$.
\end{proposition}
\begin{proof}
\sketch There is one singular simplex in every dimension and the boundary maps alternate between an isomorphism and zero.
\end{proof}

For nonempty $X$, reduced homology $\widetilde H_n(X)$ is the homology of
\[
\cdots\to C_1(X)\xrightarrow{\partial}C_0(X)\xrightarrow{\varepsilon}\mathbb{Z}\to0.
\]
Thus $H_0(X)\cong\widetilde H_0(X)\oplus\mathbb{Z}$ and $H_n(X)\cong\widetilde H_n(X)$ for $n>0$.

\section*{Homotopy Invariance}

For a map $f:X\to Y$, composition sends a singular simplex $\sigma$ to $f\circ\sigma$ and defines a chain map $f_\sharp:C_*(X)\to C_*(Y)$.

\begin{proposition}[Chain map induces homology homomorphism]
\label{prop:chain-map-induces-homology-homomorphism}
\uses{lem:boundary-composition-zero}
\dcref{ch2:2.9}
\group{Singular Homology}

A chain map induces homomorphisms on homology.
\end{proposition}
\begin{proof}
\sketch Commutation with $\partial$ carries cycles to cycles and boundaries to boundaries.
\end{proof}

The induced maps satisfy $(fg)_*=f_*g_*$ and $\mathbb{1}_*=\mathbb{1}$.

\begin{theorem}[Homotopic maps induce the same homology map]
\label{thm:homotopic-maps-induce-same-homology-map}
\uses{def:homotopy, prop:chain-map-induces-homology-homomorphism}
\dcref{ch2:2.10}
\group{Singular Homology}

If $f,g:X\to Y$ are homotopic, then $f_*=g_*:H_n(X)\to H_n(Y)$.
\end{theorem}
\begin{proof}
\sketch Subdivide $\Delta^n\times I$ into the simplices $[v_0,\ldots,v_i,w_i,\ldots,w_n]$.  The resulting prism operator $P:C_n(X)\to C_{n+1}(Y)$ satisfies $\partial P+P\partial=g_\sharp-f_\sharp$, so the two maps agree on cycles modulo boundaries.
\end{proof}

\begin{corollary}[Homotopy equivalence induces homology isomorphism]
\label{cor:homotopy-equivalence-induces-homology-isomorphism}
\uses{def:homotopy-equivalence, thm:homotopic-maps-induce-same-homology-map}
\dcref{ch2:2.11}
\group{Singular Homology}

A homotopy equivalence $f:X\to Y$ induces isomorphisms $H_n(X)\cong H_n(Y)$ for every $n$.
\end{corollary}

\begin{proposition}[Chain-homotopic maps induce the same homology map]
\label{prop:chain-homotopic-maps-induce-same-homology-map}
\uses{prop:chain-map-induces-homology-homomorphism}
\dcref{ch2:2.12}
\group{Singular Homology}

Chain-homotopic chain maps induce the same homomorphism on homology.
\end{proposition}
\begin{proof}
\sketch If $g-f=\partial P+P\partial$, then on a cycle $z$ the difference $(g-f)(z)=\partial P(z)$ is a boundary.
\end{proof}

\section*{Exact Sequences and Excision}

A sequence is exact when the image of each map equals the kernel of the next.  In particular, a short exact sequence $0\to A\to B\to C\to0$ identifies $C$ with $B/A$.

\begin{theorem}[Excision Theorem]
\label{thm:excision}
\uses{def:deformation-retraction}
\dcref{ch2:2.13}
\group{Singular Homology}

If $A\ne\varnothing$ is closed in $X$ and is a deformation retract of a neighborhood, then
\[
\cdots\to\widetilde H_n(A)\to\widetilde H_n(X)\to\widetilde H_n(X/A)\xrightarrow{\partial}\widetilde H_{n-1}(A)\to\cdots\to\widetilde H_0(X/A)\to0
\]
is exact.
\end{theorem}
\begin{proof}
\sketch Use the short exact sequence of chain complexes obtained from the inclusion of chains in $A$ and the quotient chains in $X/A$; the neighborhood deformation retraction identifies the quotient-chain homology with $\widetilde H_*(X/A)$.
\end{proof}

\begin{corollary}[Homology of spheres]
\label{cor:homology-spheres}
\uses{thm:excision}
\dcref{ch2:2.14}
\group{Singular Homology}

\[
\widetilde H_i(S^n)\cong\begin{cases}\mathbb{Z}&i=n,\\0&i\ne n.\end{cases}
\]
\end{corollary}
\begin{proof}
\sketch Apply the exact sequence to $(D^n,S^{n-1})$, use contractibility of $D^n$, and induct from $S^0$.
\end{proof}

\begin{corollary}[Non-retraction of disk onto boundary]
\label{cor:non-retraction-disk-boundary}
\uses{cor:homology-spheres}
\dcref{ch2:2.15}
\group{Singular Homology}

$\partial D^n$ is not a retract of $D^n$.  Consequently every map $D^n\to D^n$ has a fixed point.
\end{corollary}
\begin{proof}
\sketch A retraction would make the inclusion-induced map on $\widetilde H_{n-1}(S^{n-1})\cong\mathbb{Z}$ split through $\widetilde H_{n-1}(D^n)=0$.
\end{proof}

For a pair $A\subset X$, define $C_n(X,A)=C_n(X)/C_n(A)$ and
\[
H_n(X,A)=\ker(\partial:C_n(X,A)\to C_{n-1}(X,A))/\operatorname{Im}\partial.
\]
A relative cycle is a chain $\alpha$ with $\partial\alpha\in C_{n-1}(A)$; it is a relative boundary when $\alpha=\partial\beta+\gamma$ with $\gamma\in C_n(A)$.  The short exact sequence
\[
0\to C_n(A)\to C_n(X)\to C_n(X,A)\to0
\]
is a short exact sequence of chain complexes.

\begin{theorem}[The sequence of homology groups]
\label{thm:long-exact-sequence-pair}
\uses{lem:boundary-composition-zero}
\dcref{ch2:2.16}
\group{Singular Homology}

Every short exact sequence of chain complexes $0\to A\to B\to C\to0$ induces a long exact sequence
\[
\cdots\to H_n(A)\to H_n(B)\to H_n(C)\xrightarrow{\partial}H_{n-1}(A)\to H_{n-1}(B)\to\cdots.
\]
For a pair this is
\[
\cdots\to H_n(A)\to H_n(X)\to H_n(X,A)\xrightarrow{\partial}H_{n-1}(A)\to\cdots\to H_0(X,A)\to0.
\]
\end{theorem}
\begin{proof}
\sketch For a cycle $c\in C_n$ lift $c=j(b)$ and write $\partial b=i(a)$; define $\partial[c]=[a]$.  Independence of lifts and representatives follows from exactness.  Diagram chasing proves exactness at each term.
\end{proof}

For $A\ne\varnothing$, the same construction gives $\widetilde H_n(X,A)\cong H_n(X,A)$.  The connecting map sends a relative class represented by $\alpha$ to $[\partial\alpha]\in H_{n-1}(A)$.

\begin{example}[Long exact sequence of disk and sphere pair]
\label{ex:long-exact-sequence-disk-sphere}
\uses{thm:long-exact-sequence-pair}
\dcref{ch2:2.17}
\group{Singular Homology}

\[
H_i(D^n,\partial D^n)\cong\begin{cases}\mathbb{Z}&i=n,\\0&i\ne n.\end{cases}
\]
\end{example}

\begin{example}[Reduced homology isomorphic to relative homology of basepoint]
\label{ex:reduced-homology-relative}
\uses{thm:long-exact-sequence-pair}
\dcref{ch2:2.18}
\group{Singular Homology}

For $x_0\in X$, $H_n(X,x_0)\cong\widetilde H_n(X)$ for all $n$.
\end{example}

\begin{proposition}[If two maps $f,g:(X,A)\to(Y,B)$ are homotopic through maps of pairs $(X,A)\to(Y,B)$, then $f_* = g_*:H_n(X,A)\to H_n(Y,B)$]
\label{prop:homotopic-maps-pairs}
\uses{def:homotopy, thm:homotopic-maps-induce-same-homology-map}
\dcref{ch2:2.19}
\group{Singular Homology}

Homotopic maps of pairs induce the same homomorphism on relative homology.
\end{proposition}
\begin{proof}
\sketch The prism operator preserves the subcomplex of chains in $A$ and therefore descends to relative chains, where $\partial P+P\partial=g_\sharp-f_\sharp$.
\end{proof}

For $Z\subset A\subset X$, the long exact sequence of the triple is
\[
\cdots\to H_n(A,Z)\to H_n(X,Z)\to H_n(X,A)\to H_{n-1}(A,Z)\to\cdots.
\]

\begin{theorem}[Excision for relative homology]
\label{thm:excision-general}
\uses{lem:boundary-composition-zero}
\dcref{ch2:2.20}
\group{Singular Homology}

If $Z\subset A\subset X$ and $\overline Z\subset\operatorname{int}A$, then inclusion induces isomorphisms
\[
H_n(X-Z,A-Z)\cong H_n(X,A).
\]
Equivalently, if $\operatorname{int}A\cup\operatorname{int}B=X$, then $H_n(B,A\cap B)\cong H_n(X,A)$.
\end{theorem}

For an open cover $\mathcal U=\{U_j\}$, let $C_n^{\mathcal U}(X)$ be generated by singular simplices whose images lie in one $U_j$, and let $H_n^{\mathcal U}(X)$ be its homology.

\begin{proposition}[Inclusion of $\mathcal{U}$-chains is a chain homotopy equivalence]
\label{prop:inclusion-u-chains-homotopy-equivalence}
\uses{prop:chain-homotopic-maps-induce-same-homology-map}
\dcref{ch2:2.21}
\group{Singular Homology}

The inclusion $C_*^{\mathcal U}(X)\hookrightarrow C_*(X)$ is a chain homotopy equivalence and induces $H_n^{\mathcal U}(X)\cong H_n(X)$.
\end{proposition}
\begin{proof}
\sketch Barycentric subdivision repeatedly decreases simplex diameters by at most the factor $n/(n+1)$ in dimension $n$.  The induced subdivision operator $S$ is chain homotopic to the identity, and sufficiently subdivided chains are $\mathcal U$-small.
\end{proof}

\begin{proposition}[Good pairs and absolute homology]
\label{prop:good-pairs-absolute-homology}
\uses{thm:excision-general, def:deformation-retraction}
\dcref{ch2:2.22}
\group{Singular Homology}

If $(X,A)$ is a good pair, then the quotient map induces $H_n(X,A)\cong\widetilde H_n(X/A)$.
\end{proposition}
\begin{proof}
\sketch Apply excision to chains supported away from a neighborhood deformation retracting onto $A$, then identify the relative quotient with $X/A$.
\end{proof}

\begin{example}[Explicit cycles generating homology of spheres and disks]
\label{ex:cycles-spheres-disks}
\uses{cor:homology-spheres}
\dcref{ch2:2.23}
\group{Singular Homology}

For the two-simplex model of $S^n=D^n/\partial D^n$, the signed difference of the two top simplices represents the generator of $\widetilde H_n(S^n)$; the corresponding relative top cell represents the generator of $H_n(D^n,\partial D^n)$.
\end{example}

\begin{corollary}[Excision for CW subcomplexes]
\label{cor:excision-cw-subcomplexes}
\uses{thm:excision-general, def:cw-pair, prop:cw-pairs-have-hep}
\dcref{ch2:2.24}
\group{Singular Homology}

If $A$ and $B$ are CW subcomplexes and $X=A\cup B$, then inclusion induces
\[
H_n(B,A\cap B)\xrightarrow{\cong}H_n(X,A)
\]
for every $n$.
\end{corollary}

\begin{proof}
\sketch Both CW pairs are good, and the quotient spaces
$B/(A\cap B)$ and $X/A$ are naturally homeomorphic. Apply the good-pair
quotient isomorphism.
\end{proof}

\begin{corollary}[Wedge sums and direct sum of homology]
\label{cor:wedge-sum-direct-sum-homology}
\uses{prop:good-pairs-absolute-homology}
\dcref{ch2:2.25}
\group{Singular Homology}

Suppose the wedge $\bigvee_\alpha X_\alpha$ is formed at basepoints
$x_\alpha$ for which every pair $(X_\alpha,x_\alpha)$ is good. Then the
inclusions induce an isomorphism
\[
\bigoplus_\alpha\widetilde H_n(X_\alpha)
\xrightarrow{\cong}
\widetilde H_n\left(\bigvee_\alpha X_\alpha\right).
\]
\end{corollary}

\begin{proof}
\sketch Identify reduced homology with homology relative to the wedge point and
apply the good-pair quotient isomorphism before identifying the basepoints.
\end{proof}

\begin{theorem}[Invariance of dimension]
\label{thm:invariance-of-dimension}
\uses{prop:good-pairs-absolute-homology}
\dcref{ch2:2.26}
\group{Singular Homology}

If nonempty open subsets $U\subset\mathbb{R}^m$ and
$V\subset\mathbb{R}^n$ are homeomorphic, then $m=n$.
\end{theorem}
\begin{proof}
\sketch For $x\in U$, excision and the long exact sequence give
\[
H_k(U,U-\{x\})\cong H_k(\mathbb{R}^m,\mathbb{R}^m-\{x\})
\cong\widetilde H_{k-1}(S^{m-1}),
\]
which is nonzero exactly when $k=m$. A homeomorphism identifies these local
homology groups with those of $(V,V-\{h(x)\})$, which are nonzero exactly in
degree $n$; hence $m=n$.
\end{proof}

\begin{theorem}[Isomorphism of simplicial and singular homology]
\label{thm:simplicial-singular-homology-isomorphic}
\uses{lem:five-lemma}
\dcref{ch2:2.27}
\group{Singular Homology}

For every $\Delta$-complex pair $(X,A)$, the canonical map
\[
H_n^\Delta(X,A)\longrightarrow H_n(X,A)
\]
is an isomorphism for every $n$.
\end{theorem}
\begin{proof}
\sketch The characteristic maps define the canonical chain map. Compare the
exact sequences of successive skeleta; the relative groups are free on the
simplices in the relevant dimension, so induction and the five-lemma give the
absolute result for finite-dimensional $X$. Compactness reduces arbitrary
cycles and bounding chains to finite skeleta. Comparing the long exact
simplicial and singular sequences of $(X,A)$ and applying the five-lemma then
gives the relative assertion.
\end{proof}

\begin{lemma}[The Five-Lemma]
\label{lem:five-lemma}
\group{Singular Homology}

In a commutative diagram of exact sequences with five consecutive vertical maps, if the first, second, fourth, and fifth maps are isomorphisms, then so is the middle map.
\end{lemma}
\begin{proof}
\sketch Exactness transports an element in the kernel of the middle map to zero through the neighboring isomorphisms, and similarly lifts every target element.
\end{proof}

\section{Computations and Applications}

For $n>0$, the degree of $f:S^n\to S^n$ is the integer defined by $f_*[S^n]=(\deg f)[S^n]$.  It satisfies
\[
\deg\mathbb{1}=1,\quad \deg f=0\text{ if }f\text{ is not surjective},\quad f\simeq g\Rightarrow\deg f=\deg g,\quad \deg(fg)=\deg f\deg g.
\]
Reflections have degree $-1$, the antipodal map has degree $(-1)^{n+1}$, and a fixed-point-free map $S^n\to S^n$ is homotopic to the antipodal map.
Every homotopy equivalence $S^n\to S^n$ has degree $1$ or $-1$.

\begin{theorem}[Tangent vector fields on spheres]
\label{thm:tangent-vector-fields-spheres}
\uses{cor:homology-spheres}
\dcref{ch2:2.28}
\group{Cellular Homology}

$S^n$ admits a continuous nowhere-zero tangent vector field if and only if $n$ is odd.
\end{theorem}
\begin{proof}
\sketch A unit tangent field $v$ gives the homotopy $f_t(x)=\cos(t)x+\sin(t)v(x)$ from $\mathbb{1}$ to the antipodal map, forcing $(-1)^{n+1}=1$.  For $n=2k-1$, use $v(x_1,\ldots,x_{2k})=(-x_2,x_1,\ldots,-x_{2k},x_{2k-1})$.
\end{proof}

\begin{proposition}[Only Z2 acts freely on even spheres]
\label{prop:z2-acts-freely-even-spheres}
\uses{thm:tangent-vector-fields-spheres}
\dcref{ch2:2.29}
\group{Cellular Homology}

If $n$ is even, every nontrivial group acting freely on $S^n$ is isomorphic to $\mathbb{Z}_2$.
\end{proposition}
\begin{proof}
\sketch Degree defines a homomorphism from the acting group to $\{\pm1\}$.  A fixed-point-free element has degree $-1$ in even dimension, so the degree homomorphism is injective.
\end{proof}

If $f:S^n\to S^n$ has finite fiber $f^{-1}(y)=\{x_1,\ldots,x_m\}$, excision identifies the local map at $x_i$ with multiplication by the local degree $\deg f|_{x_i}$.

\begin{proposition}[Local degree formula]
\label{prop:local-degree-formula}
\uses{cor:homology-spheres}
\dcref{ch2:2.30}
\group{Cellular Homology}

\[
\deg f=\sum_{i=1}^m\deg f|_{x_i}.
\]
\end{proposition}
\begin{proof}
\sketch Excision splits $H_n(S^n,S^n-f^{-1}(y))$ into the local summands.  The global generator is the sum of the local generators, and applying $f_*$ gives the stated sum.
\end{proof}

\begin{example}[Maps of prescribed degree]
\label{ex:degree-reflection-antipodal}
\uses{cor:homology-spheres}
\dcref{ch2:2.31}
\group{Cellular Homology}

Collapsing the complement of $k$ disjoint balls in $S^n$ to a point and mapping the resulting wedge of spheres to $S^n$ produces degrees $\pm k$, by choosing the signs of the local degrees.
\end{example}

\begin{example}[Degrees of power maps and their suspensions]
\label{ex:degree-suspension}
\uses{def:suspension, cor:homology-spheres}
\dcref{ch2:2.32}
\group{Cellular Homology}

The map $z\mapsto z^k$ on $S^1$ has degree $k$, and its local degrees add to $k$.  Suspension preserves degree, so suspending this map repeatedly realizes degree $k$ on every higher sphere.
\end{example}

\begin{proposition}[Suspension preserves degree]
\label{prop:homology-cw-complexes}
\uses{def:cell-complex}
\dcref{ch2:2.33}
\group{Cellular Homology}

If $Sf:S^{n+1}\to S^{n+1}$ is the suspension of $f:S^n\to S^n$, then $\deg Sf=\deg f$.
\end{proposition}
\begin{proof}
\sketch Apply naturality to the long exact sequence of $(CS^n,S^n)$; the quotient map of pairs is $Sf$ and the boundary maps identify the two degree multiplications.
\end{proof}

\begin{lemma}[Homology of CW skeleton pairs]
\label{lem:homology-cw-skeleton-pairs}
\uses{def:cell-complex}
\dcref{ch2:2.34}
\group{Cellular Homology}

For a CW complex $X$:
\begin{enumerate}
\item $H_k(X^n,X^{n-1})=0$ for $k\ne n$, while $H_n(X^n,X^{n-1})$ is free on the $n$-cells;
\item $H_k(X^n)=0$ for $k>n$;
\item $H_k(X^n)\to H_k(X)$ is an isomorphism for $k<n$ and is surjective for $k=n$.
\end{enumerate}
Consequently, if $X$ is finite-dimensional, then $H_k(X)=0$ for
$k>\dim X$.
\end{lemma}
\begin{proof}
\sketch The quotient $X^n/X^{n-1}$ is a wedge of $n$-spheres. The pair exact
sequence gives the skeletal vanishing and stabilization maps. For arbitrary
$X$, every finite singular cycle and bounding chain has compact image and lies
in a finite skeleton, so the stabilized skeletal groups compute $H_k(X)$.
\end{proof}

The cellular chain complex is
\[
\cdots\to H_n(X^n,X^{n-1})\xrightarrow{d_n}H_{n-1}(X^{n-1},X^{n-2})\to\cdots,
\qquad d_n=j_{n-1}\partial_n.
\]
Its homology is denoted $H_n^{CW}(X)$.

\begin{theorem}[Isomorphism of cellular and singular homology]
\label{thm:cellular-singular-homology-isomorphic}
\uses{prop:homology-cw-complexes, lem:homology-cw-skeleton-pairs}
\dcref{ch2:2.35}
\group{Cellular Homology}

\[
H_n^{CW}(X)\cong H_n(X).
\]
\end{theorem}
\begin{proof}
\sketch The maps $j_n$ identify $H_n(X^n)$ with $\ker\partial_n$ and carry $\operatorname{Im}\partial_{n+1}$ onto $\operatorname{Im}d_{n+1}$, yielding the quotient isomorphism.
\end{proof}

For an $n$-cell $e_\alpha^n$, the cellular boundary is
\[
d_n(e_\alpha^n)=\sum_\beta d_{\alpha\beta}e_\beta^{n-1},
\]
where $d_{\alpha\beta}$ is the degree of the attaching map $S^{n-1}\to X^{n-1}\to X^{n-1}/(X^{n-1}-e_\beta^{n-1})\cong S^{n-1}$.

\begin{example}[Cellular homology of closed orientable surfaces]
\label{ex:cellular-homology-spheres}
\uses{thm:cellular-singular-homology-isomorphic, ex:sphere-two-cells}
\dcref{ch2:2.36}
\group{Cellular Homology}

For the genus-$g$ orientable surface with one $0$-cell, $2g$ $1$-cells, and one $2$-cell attached by $\prod_i[a_i,b_i]$, both cellular boundary maps vanish.  Thus $H_0=H_2=\mathbb{Z}$ and $H_1=\mathbb{Z}^{2g}$.
\end{example}

\begin{example}[Cellular homology of closed nonorientable surfaces]
\label{ex:cellular-homology-surfaces}
\uses{thm:cellular-singular-homology-isomorphic}
\dcref{ch2:2.37}
\group{Cellular Homology}

For the nonorientable genus-$g$ surface with attaching word $a_1^2\cdots a_g^2$, $d_2(1)=(2,\ldots,2)$ and $d_1=0$.  Hence $H_2=0$ and $H_1\cong\mathbb{Z}^{g-1}\oplus\mathbb{Z}_2$.
\end{example}

\begin{example}[An Acyclic Space]
\label{ex:acyclic-space}
\uses{thm:cellular-singular-homology-isomorphic}
\dcref{ch2:2.38}
\group{Cellular Homology}

Attaching 2-cells to $S^1\vee S^1$ by $a^5b^{-3}$ and $b^3(ab)^{-2}$ gives $d_2=\begin{pmatrix}5&-2\\-3&1\end{pmatrix}$, whose determinant is $-1$.  The resulting connected complex is acyclic, although its fundamental group is nontrivial.
\end{example}

\begin{example}[A 3-dimensional torus]
\label{ex:three-torus}
\uses{thm:cellular-singular-homology-isomorphic}
\dcref{ch2:2.39}
\group{Cellular Homology}

For $T^3$, all cellular boundary maps vanish, so $H_i(T^3)$ is $\mathbb{Z}$ for $i=0,3$, $\mathbb{Z}^3$ for $i=1,2$, and zero otherwise.  For $K\times S^1$, $d_3(c^3)=2C$, $d_2(A)=2b$, and $d_2(B)=d_2(C)=0$, giving $H_3=0$, $H_2\cong\mathbb{Z}\oplus\mathbb{Z}_2$, and $H_1\cong\mathbb{Z}^2\oplus\mathbb{Z}_2$.
\end{example}

\begin{example}[Moore Spaces]
\label{ex:moore-spaces}
\uses{thm:cellular-singular-homology-isomorphic}
\dcref{ch2:2.40}
\group{Cellular Homology}

A Moore space $M(G,n)$ is a simply-connected CW complex for $n>1$ with $H_n(M(G,n))\cong G$ and reduced homology zero in other dimensions.  For finitely generated $G$, wedge spheres and degree maps realize it; in general attach $(n+1)$-cells representing a free presentation of $G$.
\end{example}

\begin{example}[CW complexes realizing prescribed homology groups]
\label{ex:cellular-homology-cpn}
\uses{thm:cellular-singular-homology-isomorphic, ex:complex-projective-space-cw}
\dcref{ch2:2.41}
\group{Cellular Homology}

Wedge sums of Moore spaces realize any prescribed sequence of homology groups in positive dimensions.
\end{example}

\begin{example}[Cellular homology of $\mathbb{RP}^n$]
\label{ex:cellular-homology-rpn}
\uses{thm:cellular-singular-homology-isomorphic, ex:real-projective-space-cw}
\dcref{ch2:2.42}
\group{Cellular Homology}

$\mathbb{RP}^n$ has one cell in each dimension $0\leq k\leq n$, with cellular boundary $d_k=0$ for odd $k$ and multiplication by $2$ for even $k$.  Hence
\[
H_k(\mathbb{RP}^n)=\begin{cases}\mathbb{Z}&k=0,\text{ or }k=n\text{ odd},\\\mathbb{Z}_2&0<k<n\text{ odd},\\0&\text{otherwise}.\end{cases}
\]
\end{example}

\begin{example}[Lens Spaces]
\label{ex:lens-spaces}
\uses{thm:cellular-singular-homology-isomorphic}
\dcref{ch2:2.43}
\group{Cellular Homology}

For $m>1$ and $(\ell_j,m)=1$, let $L_m(\ell_1,\ldots,\ell_n)=S^{2n-1}/\mathbb{Z}_m$ under $(z_j)\mapsto(e^{2\pi i\ell_j/m}z_j)$.  Its cellular chain complex has one copy of $\mathbb{Z}$ in each dimension through $2n-1$, with boundary maps alternating between $0$ and multiplication by $m$.  Therefore
\[
H_k(L_m(\ell_1,\ldots,\ell_n))=\begin{cases}\mathbb{Z}&k=0,2n-1,\\\mathbb{Z}_m&0<k<2n-1\text{ odd},\\0&\text{otherwise}.\end{cases}
\]
\end{example}

\begin{theorem}[Euler characteristic in terms of homology]
\label{thm:euler-characteristic-homology}
\uses{def:dimension-cw-complex}
\dcref{ch2:2.44}
\group{Cellular Homology}

For a finite CW complex,
\[
\chi(X)=\sum_n(-1)^n\#\{n\text{-cells}\}=\sum_n(-1)^n\operatorname{rank}H_n(X).
\]
\end{theorem}
\begin{proof}
\sketch In a finite chain complex, use the short exact sequences $0\to Z_n\to C_n\to B_{n-1}\to0$ and $0\to B_n\to Z_n\to H_n\to0$, add the resulting rank identities with alternating signs, and apply this to the cellular complex.
\end{proof}

\begin{lemma}[Splitting Lemma]
\label{lem:splitting-lemma}
\group{Singular Homology}

For $0\to A\xrightarrow{i}B\xrightarrow{j}C\to0$, the following are equivalent: there is $p:B\to A$ with $pi=\mathbb{1}$; there is $s:C\to B$ with $js=\mathbb{1}$; and $B\cong A\oplus C$ compatibly with the sequence.
\end{lemma}
\begin{proof}
\sketch A retraction gives $b\mapsto(p(b),j(b))$; a section gives $(a,c)\mapsto i(a)+s(c)$.  These maps are inverse isomorphisms, and the converse maps are projections and inclusions.
\end{proof}

\begin{proposition}[Fundamental group of finite-dimensional K(G,1) is torsionfree]
\label{prop:fundamental-group-kg1-torsionfree}
\uses{def:k-g-1-space, ex:lens-spaces}
\dcref{ch2:2.45}
\group{Cellular Homology}

If a finite-dimensional CW complex $X$ is a $K(G,1)$, then $G=\pi_1(X)$ is torsionfree.
\end{proposition}
\begin{proof}
\sketch A finite cyclic subgroup would give a finite-dimensional covering $K(\mathbb{Z}_m,1)$, but $H_{2r+1}(\mathbb{Z}_m)\cong\mathbb{Z}_m$ for all $r\geq0$, contradicting finite dimensionality.
\end{proof}

\section*{Mayer--Vietoris Sequences}

If $X=A\cup B$ with $\operatorname{int}A\cup\operatorname{int}B=X$, then
\[
\cdots\to H_n(A\cap B)\xrightarrow{\Phi}H_n(A)\oplus H_n(B)\xrightarrow{\psi}H_n(X)\xrightarrow{\delta}H_{n-1}(A\cap B)\to\cdots.
\]
At chain level use $\Phi(z)=(z,-z)$ and $\psi(x,y)=x+y$; the small-chain proposition identifies the subcomplex of sums of chains in $A$ and $B$ with $C_*(X)$.

\begin{example}[Mayer-Vietoris sequence of sphere split into hemispheres]
\label{ex:mayer-vietoris-sphere}
\uses{thm:excision-general}
\dcref{ch2:2.46}
\group{Singular Homology}

For the hemispherical decomposition $S^n=A\cup B$ with $A\cap B=S^{n-1}$, $\widetilde H_i(S^n)\cong\widetilde H_{i-1}(S^{n-1})$.
\end{example}

\begin{example}[Mayer-Vietoris sequence of Klein bottle]
\label{ex:mayer-vietoris-klein-bottle}
\uses{thm:excision-general}
\dcref{ch2:2.47}
\group{Singular Homology}

Writing the Klein bottle as two Möbius bands gives $\Phi:\mathbb{Z}\to\mathbb{Z}^2$, $1\mapsto(2,-2)$.  Thus $H_2(K)=0$ and $H_1(K)\cong\mathbb{Z}\oplus\mathbb{Z}_2$.
\end{example}

\begin{example}[Homology of mapping torus]
\label{ex:homology-mapping-torus}
\uses{thm:excision-general}
\dcref{ch2:2.48}
\group{Singular Homology}

For maps $f,g:X\to Y$, let $Z=(X\times I\sqcup Y)/((x,0)\sim f(x),(x,1)\sim g(x))$.  Then
\[
\cdots\to H_n(X)\xrightarrow{f_*-g_*}H_n(Y)\to H_n(Z)\to H_{n-1}(X)\xrightarrow{f_*-g_*}H_{n-1}(Y)\to\cdots.
\]
For the mapping torus of a reflection of $S^1$, this gives $H_2=0$ and $H_1\cong\mathbb{Z}\oplus\mathbb{Z}_2$.
\end{example}

\textbf{Homology with Coefficients}

For an abelian group $G$, define $C_n(X;G)$ by allowing coefficients in $G$, set $C_n(X,A;G)=C_n(X;G)/C_n(A;G)$, and use the alternating boundary formula.  The resulting groups are $H_n(X;G)$ and $H_n(X,A;G)$; reduced groups use $C_0(X;G)\xrightarrow{\varepsilon}G$.  All chain, homotopy, exact-sequence, excision, and Mayer--Vietoris constructions extend to these coefficients.  In particular, $\widetilde H_i(S^k;G)$ is $G$ for $i=k$ and zero otherwise, and cellular chains have one copy of $G$ per cell.

\begin{lemma}[Homology map of degree m map with G coefficients]
\label{lem:homology-map-degree-m-coefficients}
\uses{cor:homology-spheres}
\dcref{ch2:2.49}
\group{Singular Homology}

If $f:S^k\to S^k$ has degree $m$, then $f_*:H_k(S^k;G)\to H_k(S^k;G)$ is multiplication by $m$.
\end{lemma}
\begin{proof}
\sketch Naturality with respect to the coefficient map $\mathbb{Z}\to G$, $1\mapsto g$, transfers the degree-$m$ action on integral sphere homology to $g\mapsto mg$.
\end{proof}

\begin{example}[Homology of RPn with field coefficients]
\label{ex:homology-rpn-field-coefficients}
\uses{ex:cellular-homology-rpn}
\dcref{ch2:2.50}
\group{Singular Homology}

The cellular complex of $\mathbb{RP}^n$ over a field $F$ alternates $0$ and multiplication by $2$.  If $\operatorname{char}F=2$, then $H_k(\mathbb{RP}^n;F)\cong F$ for $0\leq k\leq n$; otherwise it is $F$ for $k=0$ and for $k=n$ odd, and zero otherwise.
\end{example}

\begin{example}[Homology of Moore space quotient]
\label{ex:homology-moore-space-quotient}
\uses{ex:moore-spaces}
\dcref{ch2:2.51}
\group{Singular Homology}

For $X=M(\mathbb{Z}_m,n)$ and $f:X\to X/S^n=S^{n+1}$, $f_*$ is zero on
reduced integral homology because the nonzero groups of source and target occur
in different dimensions. With $\mathbb{Z}_m$ coefficients, the cellular
differential $\mathbb{Z}_m\xrightarrow{m}\mathbb{Z}_m$ is zero, so
$\widetilde H_{n+1}(X;\mathbb{Z}_m)\cong\mathbb{Z}_m$. The long exact sequence
of $(X,S^n)$ begins with zero and therefore makes
\[
f_*:\widetilde H_{n+1}(X;\mathbb{Z}_m)
\longrightarrow\widetilde H_{n+1}(S^{n+1};\mathbb{Z}_m)
\]
injective, hence nonzero.
\end{example}

\section{The Formal Viewpoint}\label{sec:formal-viewpoint}

A reduced homology theory on nonempty CW complexes assigns groups $\widetilde h_n(X)$ and induced maps satisfying identity and composition laws, homotopy invariance, a natural long exact sequence for CW pairs, and the wedge axiom
\[
\bigoplus_\alpha\widetilde h_n(X_\alpha)\xrightarrow{\cong}\widetilde h_n\left(\bigvee_\alpha X_\alpha\right).
\]
Unreduced theories instead use relative groups and an excision axiom.  Reduced and unreduced theories are related by $\widetilde h_n(X)=\ker(h_n(X)\to h_n(\mathrm{point}))$ and $h_n(X)=\widetilde h_n(X_+)$.

For $x_0\in X$, the split exact sequence of $(X,x_0)$ gives
\[
h_n(X)\cong\widetilde h_n(X)\oplus h_n(x_0).
\]
The groups $h_n(x_0)\cong\widetilde h_n(S^0)$ are the coefficient groups.
Any sequence of abelian groups $G_i$ occurs as coefficients of the theory
$h_n(X,A)=\bigoplus_i H_{n-i}(X,A;G_i)$. Although coefficients do not determine
an arbitrary homology theory, a theory on CW pairs satisfying the dimension
axiom is naturally singular homology with coefficients in $G=h_0(x_0)$.

A category consists of objects, morphism sets with identity morphisms, and associative composition.  A functor preserves identities and composition; a contravariant functor reverses arrows.  A natural transformation $T:F\Rightarrow G$ is a family of morphisms $T_X:F(X)\to G(X)$ commuting with every morphism.  Singular chains and homology are covariant functors, while cohomology is contravariant.  The classifying space $B\mathcal C$ of a category has simplices given by strings of composable morphisms; functors induce maps and natural transformations induce homotopies.

\section{2.A Homology and Fundamental Group}

\begin{theorem}[Hurewicz theorem in dimension 1]
\label{thm:hurewicz-dimension-1}
\uses{def:loop-and-pi1-set, prop:h0-path-connected}
\dcref{ch2:2A.1}
\group{Singular Homology}

For path-connected $X$, regarding based loops as singular 1-cycles gives a surjection
\[
h:\pi_1(X,x_0)\twoheadrightarrow H_1(X)
\]
whose kernel is the commutator subgroup.  Hence $H_1(X)\cong\pi_1(X,x_0)_{\mathrm{ab}}$.
\end{theorem}
\begin{proof}
\sketch Constant paths bound constant 2-simplices, homotopies give prism boundaries, concatenation becomes addition, and reversal becomes negation.  These facts define $h$.  A cycle can be joined to the basepoint and concatenated into one loop, proving surjectivity.  If a loop bounds a 2-chain, pair the boundary edges of its triangles to form a $\Delta$-complex; in the abelianized fundamental group each triangle boundary is trivial, so the loop lies in the commutator subgroup.
\end{proof}

\begin{example}[H1 of closed orientable surface]
\label{ex:h1-closed-orientable-surface}
\uses{thm:hurewicz-dimension-1}
\dcref{ch2:2A.2}
\group{Singular Homology}

For a closed orientable genus-$g$ surface $M$, $H_1(M)\cong\mathbb{Z}^{2g}$, generated by the standard $\alpha_i,\beta_i$ cycles.
\end{example}

\section{Classical Applications}

\begin{proposition}[Homology of sphere complement of disk/sphere]
\label{prop:homology-sphere-complement-disk-sphere}
\uses{thm:excision-general, ex:mayer-vietoris-sphere}
\dcref{ch2:2B.1}
\group{Singular Homology}

If $h:D^k\hookrightarrow S^n$, then $\widetilde H_i(S^n-h(D^k))=0$ for every $i$.  If $h:S^k\hookrightarrow S^n$ with $k<n$, then
\[
\widetilde H_i(S^n-h(S^k))\cong\begin{cases}\mathbb{Z}&i=n-k-1,\\0&\text{otherwise}.\end{cases}
\]
\end{proposition}
\begin{proof}
\sketch For embedded disks, induct on $k$ after writing $D^k=I^k$. Splitting
the last coordinate in half and applying Mayer--Vietoris shows that a
nonbounding cycle in the complement would remain nonbounding in one half.
Iteration produces nested intervals shrinking to a point. By induction the
cycle bounds outside the corresponding embedded $(k-1)$-disk; compactness of
the finite bounding chain places it outside one sufficiently small interval,
a contradiction. Thus every disk complement has trivial reduced homology.
For an embedded $S^k$, split it into two hemispheres. Their complements have
trivial reduced homology by the disk case, and Mayer--Vietoris yields
$\widetilde H_i(S^n-h(S^k))\cong
\widetilde H_{i+1}(S^n-h(S^{k-1}))$. Induction from $S^0$ gives the formula.
\end{proof}

\begin{example}[The Alexander Horned Sphere]
\label{ex:alexander-horned-sphere}
\uses{prop:homology-sphere-complement-disk-sphere}
\dcref{ch2:2B.2}
\group{Singular Homology}

There is an embedding $S^2\hookrightarrow\mathbb{R}^3$ whose unbounded complement is not simply-connected.  The nested solid-torus construction has complements with free groups $F_{2^n}$, where each old generator becomes a commutator of two new generators; their union is the complement group and has trivial abelianization.
\end{example}

\begin{theorem}[Invariance of domain]
\label{thm:invariance-of-domain-general}
\uses{prop:homology-sphere-complement-disk-sphere}
\dcref{ch2:2B.3}
\group{Singular Homology}

If $U\subset\mathbb{R}^n$ is open and $h:U\to\mathbb{R}^n$ is a continuous injection, then $h(U)$ is open.
\end{theorem}
\begin{proof}
\sketch Compactify to $S^n$.  The complement of the embedded boundary of a small disk has two components by the complement proposition; one is the image of the disk interior and is therefore open.
\end{proof}

\begin{corollary}[Embedding of compact manifold in connected manifold]
\label{cor:embedding-compact-manifold-connected}
\uses{thm:invariance-of-domain-general}
\dcref{ch2:2B.4}
\group{Singular Homology}

An embedding of a compact $n$-manifold into a connected $n$-manifold is surjective and hence a homeomorphism.
\end{corollary}
\begin{proof}
\sketch The image is closed by compactness and open by invariance of domain; connectedness forces it to be all of the target.
\end{proof}

\begin{theorem}[Real and complex division algebras]
\label{thm:real-complex-division-algebras}
\uses{prop:z2-acts-freely-even-spheres}
\dcref{ch2:2B.5}
\group{Cellular Homology}

$\mathbb{R}$ and $\mathbb{C}$ are the only finite-dimensional commutative division algebras over $\mathbb{R}$ with identity.
\end{theorem}
\begin{proof}
\sketch In a commutative division algebra $A=\mathbb{R}^n$, $x\mapsto x^2/|x|^2$ descends to an injective map $\mathbb{RP}^{n-1}\to S^{n-1}$.  Invariance of domain makes it onto, and homology forces $n=1$ or $2$.  The identity and a nonreal element with square $-1$ identify the two-dimensional case with $\mathbb{C}$.
\end{proof}

\begin{proposition}[Odd map of Sn has odd degree]
\label{prop:odd-map-odd-degree}
\uses{cor:homology-spheres}
\dcref{ch2:2B.6}
\group{Cellular Homology}

If $f:S^n\to S^n$ satisfies $f(-x)=-f(x)$, then $\deg f$ is odd.
\end{proposition}
\begin{proof}
\sketch Use the transfer exact sequence for $S^n\to\mathbb{RP}^n$ with $\mathbb{Z}_2$ coefficients.  An odd map induces a morphism of transfer sequences; induction and exactness show that $f_*$ on top homology is an isomorphism.  The degree map with these coefficients is multiplication by $\deg f\bmod2$.
\end{proof}

\begin{corollary}[Borsuk-Ulam theorem]
\label{cor:borsuk-ulam}
\uses{prop:odd-map-odd-degree}
\dcref{ch2:2B.7}
\group{Cellular Homology}

For every $g:S^n\to\mathbb{R}^n$, some $x$ satisfies $g(x)=g(-x)$.
\end{corollary}
\begin{proof}
\sketch Otherwise $g(x)-g(-x)$ normalizes to an odd map $S^n\to S^{n-1}$.  Its restriction to the equator has odd degree but extends over a hemisphere and is nullhomotopic, a contradiction.
\end{proof}

\section{Simplicial Approximation}

A simplicial map sends vertices to vertices and is affine on every simplex.  A vertex map extends to a simplicial map exactly when the vertices of every source simplex map into a target simplex.  For a simplex $\sigma$, write $\operatorname{st}\sigma$ for the union of interiors of simplices containing $\sigma$.

\begin{theorem}[Simplicial Approximation Theorem]
\label{thm:simplicial-approximation}
\uses{lem:star-intersections}
\dcref{ch2:2C.1}
\group{Singular Homology}

If $K$ is finite and $L$ is any simplicial complex, every continuous $f:K\to L$ is homotopic to a map simplicial on some iterated barycentric subdivision of $K$.
\end{theorem}
\begin{proof}
\sketch Choose a Lebesgue number for the inverse images of vertex stars.  Subdivide $K$ until closed vertex stars have smaller diameter than this number, choose for each source vertex a target vertex whose star contains its image, and extend linearly.  The star-intersection lemma makes the extension simplicial; the two maps lie in common target simplices, so straight-line interpolation gives the homotopy.
\end{proof}

\begin{lemma}[Intersections of star neighborhoods]
\label{lem:star-intersections}
\uses{def:simplex}
\dcref{ch2:2C.2}
\group{Singular Homology}

For vertices $v_1,\ldots,v_n$ of a simplicial complex, $\operatorname{st}v_1\cap\cdots\cap\operatorname{st}v_n$ is empty unless they span a simplex $\sigma$, in which case the intersection equals $\operatorname{st}\sigma$.
\end{lemma}
\begin{proof}
\sketch A point in the intersection lies in the interior of a simplex containing all the $v_i$; the simplices containing those vertices are precisely the cofaces of $\sigma$.
\end{proof}

\begin{theorem}[Lefschetz Fixed Point Theorem]
\label{thm:lefschetz-fixed-point}
\uses{thm:cellular-singular-homology-isomorphic}
\dcref{ch2:2C.3}
\group{Singular Homology}

For a finite simplicial complex, or a retract of one, if $\tau(f)=\sum_n(-1)^n\operatorname{tr}(f_*:H_n(X)\to H_n(X))\ne0$, then $f$ has a fixed point.
\end{theorem}
\begin{proof}
\sketch Reduce from a retract to the finite complex.  If a finite-complex map has no fixed point, approximate it after subdivision by a simplicial map $g$ with $g(\sigma)\cap\sigma=\varnothing$ for every simplex.  The induced cellular matrices have zero diagonal, so the alternating trace is zero.  Homotopy invariance gives $\tau(f)=\tau(g)=0$.
\end{proof}

\begin{example}[Lefschetz number of surface rotation]
\label{ex:lefschetz-surface-rotation}
\uses{thm:lefschetz-fixed-point}
\dcref{ch2:2C.4}
\group{Singular Homology}

For the fixed-point-free half-turn of a genus-three orientable surface, the induced maps contribute $1$ in dimension $0$, $-2$ in dimension $1$, and $1$ in dimension $2$, so $\tau=0$.
\end{example}

\section{Simplicial Approximations to CW Complexes}

\begin{theorem}[CW complex homotopy equivalent to simplicial complex]
\label{thm:cw-homotopy-equivalent-simplicial}
\uses{def:cell-complex}
\dcref{ch2:2C.5}
\group{CW Complexes}

Every CW complex is homotopy equivalent to a simplicial complex of the same dimension; finiteness and countability are preserved.
\end{theorem}
\begin{proof}
\sketch Inductively replace the attaching maps of the skeleta by simplicial maps.  For a simplicial attaching map $f$, construct a simplicial mapping cylinder containing a subdivision of the domain and retracting onto the target, then attach a cone to obtain a simplicial mapping cone.  The resulting simplicial complex $Y_{n+1}$ and the original skeleton $X^{n+1}$ are deformation retracts of a common complex.  Taking unions gives $Y\simeq X$.
\end{proof}
